\section{\label{sec:design}Design}
Let $\theta_0$ be a small instruction-tuned open-weight LLM.
Each run applies a privacy mechanism $m_i$ at a given strength,
acting on either the adaptation data or the adaptation procedure,
and yields two measurements.
\textbf{Semantic distortion} $S_i$ is the mean cosine distance between sentence embeddings of the original and the rewritten training text;
where the mechanism leaves the text intact,
it is measured between generations of the adapted and the base model on a fixed probe set,
and in both cases it is computed independently of the privacy parameter that produced it.
\textbf{Alignment degradation} $Y_i$ is the drop from the same model's non-private baseline on the suite of \S\ref{sec:eval}.

The two candidate pathways make opposite predictions about what survives conditioning on distortion.
If degradation runs through meaning,
the association between privacy level and $Y$ vanishes once $S$ is held fixed;
if it travels the other route,
a level effect remains at fixed $S$. Splitting the association that way and writing it linearly yields
\begin{equation}
  Y_i = a_{m_i} + b_{m_i} d_i + c S_i + e_i,
  \label{eq:model}
\end{equation}
where $d_i$ is the privacy-level rank of run $i$ within its own mechanism,
$a_m$ and $b_m$ are a per-mechanism intercept and level slope,
and $e_i$ is the error.
A random intercept per base model is added to $a_m$;
$Y$,
$d$ and $S$ are z-scored,
so $b_m$ and $c$ are standardized effects on one scale.
Baseline runs have $Y \equiv 0$ by definition and stand outside the regression.
Despite the mechanism index,
\eqref{eq:model} is one regression estimated jointly through indicator variables,
so each row activates the intercept and slope of its own mechanism only while the $S$ column stays shared across all of them.
That makes $c$ shared,
pooling mechanisms that disagree about how much distortion a given privacy level buys,
and puts $\hat c$ and every $\hat b_m$ in one covariance matrix,
so comparing them is a Wald test.
The level enters only through its rank,
because the nominal budgets are incomparable across families.

We state the claim in two stages,
because a shared coefficient can be distinct from zero and still too small to matter.
The first fixes a negligibility floor $\Delta = 0.10$ in standardized units,
committed before the main pass together with the primary $S$ metric and the exclusion rules,
and asks whether $|c| \ge \Delta$;
an interval for $c$ inside $(-\Delta, \Delta)$ would settle it the other way by a two one-sided test (TOST)~\citep{schuirmann1987comparison, lakens2018equivalence}.
The second is substantive:
$|c| > |b_m|$ for every mechanism,
distortion predicting degradation more strongly than the nominal level does.

Within a single mechanism $S$ and $\varepsilon$ are correlated by construction,
so separating $b_m$ from $c$ takes variation across mechanisms.
Variation in $S$ at fixed privacy comes from two sources.
The first is the contrast between families,
which reach comparable privacy with different amounts of semantic damage,
since sanitization passes out-of-vocabulary tokens through untouched while random replacement destroys them and identifying information is typically out of vocabulary (\S\ref{sec:mech}).
The second is the non-private control:
if it degrades alignment as much as the private conditions at equal $S$,
semantic damage rather than privacy is the cause;
if it falls short,
the residual measures mechanism-specific,
non-semantic effects.
We report variance inflation factors (VIF) for $S$,
falling back to the within-level cross-mechanism contrast above ten.
\apx{extended} carries the full development.
